\section{Related Work}
\label{sec:related}

\paragraph{RCA for microservices.}
The dominant body of work on microservice reliability targets root-cause analysis. Fu et
al.~\cite{fu2025} survey the field and identify a persistent gap between benchmark performance and
production deployability, driven largely by false positives. Soldani and Brogi~\cite{soldani2022}
catalogue microservice anti-patterns and failure modes; we reuse their taxonomy, together with
Gregg's USE methodology~\cite{gregg2013} and cascading-failure modelling~\cite{aniello2014}, to
anchor our OWL class hierarchy in the literature rather than in ad-hoc cluster labels. EWAT
departs from this line of work in objective: we predict and type developing faults rather than
explaining past ones.

\paragraph{Concept drift.}
Distinguishing benign distributional change from genuine anomalies is a concept-drift problem.
Hinder et al.~\cite{hinder2024} formalise drift in unsupervised streams, and Myrtollari et
al.~\cite{myrtollari2025} propose drift-aware anomaly detection specifically for Kubernetes
microservices. EWAT operationalises the drift/anomaly separation as an explicit fourth regime
$\thetadn$ and a look-through filter built on the kernel two-sample test~\cite{gretton2012},
estimated cheaply with Random Fourier Features. Our negative result on look-through
(Section~\ref{sec:validity}) is, to our knowledge, a useful data point on the limits of MMD-based
separation on short episodes.

\paragraph{Representation learning on service graphs.}
We encode the time-evolving service graph with a spatio-temporal graph convolutional network
(STGCN)~\cite{yu2018}, and compare against graph attention~\cite{velivckovic2018} and contrastive
pre-training in the SimCLR style~\cite{chen2020simclr}. Type structure is assessed with the
silhouette criterion, using the conventional acceptability threshold of $0.3$~\cite{kaufman1990},
complemented by the gap statistic~\cite{tibshirani2001} for choosing the number of clusters.

\paragraph{Causality and interpretability.}
For the ontology we estimate directed information flow with Transfer Entropy using the KSG
estimator~\cite{kraskov2004} --- never Granger causality, which assumes linearity inappropriate
for our signals --- and control multiple testing with Benjamini--Hochberg FDR~\cite{benjamini1995}.
Feature attributions are validated with SHAP~\cite{lundberg2017} after we found gradient-based
saliency to be unreliable ($\rho=-0.34$ against permutation importance).

\paragraph{Open-set recognition.}
Because production faults include types never seen in training, we evaluate open-set recognition
with OpenMax~\cite{bendale2016openmax} and discuss Mahalanobis~\cite{lee2018mahalanobis} and
energy-based~\cite{liu2020energy} alternatives as future work. For external validation we adapt
the public RCAEval RE2-OB benchmark~\cite{rcaeval}.
